\section{Discussion}
\label{sec:discussion}

\begin{reviewblueblock}
\noindent\textbf{What the analysis explains.}~
Section~\ref{sec:formulation} separates four effects of intermittent summaries: drift of represented classes, loss of class coverage, the reference gap introduced by learned weighting, and the representation error of Gaussian synthesis. Only the first enters Theorem~\ref{thm:convergence}. Propositions~\ref{prop:replay_drift}--\ref{prop:weighting_shift} identify which design parameters control each effect: the mixing coefficient $\alpha$ and the retention limit for replay, the mask rate and the number of sources holding a class for masking, and the disagreement between sources for weighting. These results are structural. They do not show that \ehsfp{} improves on \hsfp{}; they show where each mechanism can help or hurt.

\noindent\textbf{What the experiments show.}~
In Table~\ref{tab:ablation}, learned weighting has the highest mean accuracy among the aggregation rows labeled $\memcap=0$. The reported full-synthesis configuration exceeds that row by 0.39 percentage points in the reference condition, 1.35 under source masking, and 1.61 under delay. These are descriptive contrasts because the run-specific overrides are unavailable. The reference-condition difference is smaller than the five-run standard deviations, the masking and delay differences exceed them only modestly, and no significance tests are reported. The comparison with legacy uniform aggregation also changes how the variance is computed. The baseline rows in Table~\ref{tab:main_classification} come from a different protocol. Measuring drift and coverage directly in simulation, where fresh references are available, would provide a stronger test of the analysis than accuracy alone.

\noindent\textbf{Limitations.}~
Severe statistical heterogeneity remains the main limitation. Under $\mathrm{Dir}(0.30)$, \ehsfp{} reaches 9.13\% accuracy on CIFAR-100 and 33.08\% on HAM10000, well below the corresponding IID scores. Remark~\ref{rem:scatter_pooling} and Proposition~\ref{prop:weighting_shift} suggest two contributing factors: the pooled covariance omits between-source variation, and aggregated means become sensitive to weighting when sources disagree. Unimodal Gaussian synthesis also cannot represent multimodal class-conditional features (Remark~\ref{rem:representation_error}). On the theoretical side, the bound holds for one upper-tier learner under an SGD-type update model, whereas the experiments use Adam. The residual terms need not vanish, and the age-based bound requires pipeline sensitivity and a finite age bound on every contribution, including contributions that an untruncated accumulator retains after record expiry. No result links tier-wise stationarity to the composed prediction risk~\eqref{eq:global_objective}.

\noindent\textbf{Deployment.}~
The Adam experiments characterize task performance in round-driven simulations with persistent learning state. Deployment measurements would need to account for scatter transfers and factorization when full covariance is used, image-specific feature and mask transfers during ISIC-2018 decoder training, and separate model synchronization. These operations affect total communication and may influence latency, throughput, and cost.
\end{reviewblueblock}

